% Korean Noether P09 fresh completion tail — producer draft, unchecked.
% Current German authority: ED0004 lines 7330--7680.
% Authority SHA-256: 0CB422ECD397DD392A8625297A508DAEE3A5A934EA19EEEF49B47B319EA4F2BB.
% Translation only; independent checking and certification remain external.

\section*{§ 8. 영역 $\mathfrak G$와 $\mathfrak H$의 유별 분류.\srcfn{*)}{E. Steinitz, \emph{Algebraische Theorie der Körper}, Crelles Journal 137 (1910), 특히 § 23과 § 24 참조. (§ 23에서는 이미 초월기저의 존재가 증명되어 있다.)}}

가장 일반적인 영역 $\mathfrak G$와 $\mathfrak H$에 관한 물음과 더불어, 아래에서 더 정확히 규정할 의미에서 같은 구성 원리로 만들어 낼 수 없는 \emph{본질적으로 서로 다른} 영역들에 관한 물음이 제기된다. 이는 이 영역들을 유별로 분류하는 데 이른다.

\emph{두 영역의 원소들을 일대일로 대응시켜, 한 영역의 임의의 두 원소의 합과 곱에 언제나 대응하는 원소들의 합과 곱이 대응하게 할 수 있을 때, 두 영역은 같은 유별에 속한다거나 서로 동형이라고 하기로 한다.} 이 정의에 따르면 모든 동형 대응에서 영과 단위원은 각각 자기 자신에 대응하고, 유한개의 원소 사이에 성립하는 모든 대수적 관계는 대응하는 원소들에 대해서도 보존되며 그 역도 성립한다. 따라서 단위원뿐 아니라 모든 유리수도 자기 자신에 대응한다. 반면 대수적 수 전체와 마찬가지로 대수적 정수 전체는 각각 자기 자신 안으로 옮겨지지만, 개별 대수적 수에는 그 켤레가 대응할 수도 있다.

먼저, 고정된 초월기저 $H$를 포함하는 모든 영역 $\mathfrak G$의 전체 $\mathfrak M(\mathfrak G)$에 속하는 각 영역은 다른 초월기저 $H^*$에 딸린 전체 $\mathfrak M^*(\mathfrak G)$의 한 영역과 동형임을 지적해 두자. 실제로 모든 복소수의 체는 어느 초월기저에서든 대수적 확대를 거쳐 생기므로 그 기저와 같은 기수를 가진다.\srcfn{**)}{Steinitz, § 23과 § 24.} 따라서 $H$와 $H^*$는 같은 기수를 가지며, 기저원소 $\eta_i$와 $\eta_i^*$를 서로 대응시키면 원하는 동형사상을 얻는다.\srcfn{***)}{물론 이 논법은 두 유리기저 $\Theta$와 $\Theta^*$에 딸린 집합 $\mathfrak N(\mathfrak G)$와 $\mathfrak N^*(\mathfrak G)$에는 적용되지 않는다. $\Theta$와 $\Theta^*$를 서로 유리적으로 나타낼 수 있다는 사실만으로는 $\mathfrak N(\mathfrak G)$와 $\mathfrak N^*(\mathfrak G)$가 동형이라고 결론지을 수 없다. 다만 이 사실은 체 $(\Theta)$와 $(\Theta^*)$ 사이의 동형 대응을 준다.} 그러나 § 2, § 3, § 5의 예에서 알 수 있듯이 $\mathfrak M(\mathfrak G)$는 서로 동형이 아닌, 곧 본질적으로 서로 다른 영역 $\mathfrak G$들을 포함한다. 동형사상은 모든 대수적 관계를 보존하므로 초월기저에는 다시 초월기저가 대응해야 한다. 그런데 § 2와 § 3에서는 함수식 영역과 모듈 영역에 체르멜로 영역의 모든 기저 성질을 만족하는 초월기저가 존재하지 않음을 보였다. § 2의 함수식 영역은 대수적 정수성에 관하여 닫혀 있으므로, 부분집합 $\mathfrak M(\mathfrak H)$도 본질적으로 서로 다른 영역 $\mathfrak H$들을 포함한다.

이제 기저를 구성할 때와 마찬가지로 $\mathfrak M(\mathfrak G)$에서 다음 성질을 가진 부분집합 $\mathfrak L(\mathfrak G)$를 골라낼 수 있음을 보이겠다. $\mathfrak L(\mathfrak G)$에 속하는 영역 $\mathfrak G$들은 모두 본질적으로 서로 다르지만, $\mathfrak M(\mathfrak G)$의 각 영역, 따라서 모든 영역 $\mathfrak G$는 $\mathfrak L(\mathfrak G)$의 한 영역과 동형이다. $\mathfrak M(\mathfrak G)$에 정렬순서 $\Omega$를 부여하자. 그러면 $\mathfrak M(\mathfrak G)$의 영역 $\mathfrak G$는 앞선 어떤 영역과 동형일 수도 있고 그렇지 않을 수도 있다. 앞선 어느 영역과도 동형이 아닌 영역 $\mathfrak G$들은 집합 $\mathfrak L(\mathfrak G)$를 이룬다. 구성상 이 집합에는 본질적으로 서로 다른 영역 $\mathfrak G$만 들어 있다. 또한 귀납 논법에 따르면 $\mathfrak M(\mathfrak G)$의 각 영역은 $\mathfrak L(\mathfrak G)$의 한 영역과 동형이다. 실제로 이 성질을 만족하지 않는 첫 영역이 $\mathfrak G_1$이라고 하자. $\mathfrak G_1$은 $\mathfrak L(\mathfrak G)$에 속하거나, 앞선 영역 $\mathfrak G_0$와 동형이다. 후자의 경우 가정에 따라 $\mathfrak G_0$가 $\mathfrak L(\mathfrak G)$의 한 영역과 동형이므로 $\mathfrak G_1$도 그러하다. 모든 영역 $\mathfrak G$는 어떤 집합 $\mathfrak M^*(\mathfrak G)$에 속하고, 따라서 $\mathfrak M(\mathfrak G)$의 한 영역과 동형이다. \emph{그러므로 $\mathfrak L(\mathfrak G)$는 영역 $\mathfrak G$의 모든 유별을 대표한다.}

$\mathfrak L(\mathfrak G)$의 영역 $\mathfrak G$에서 기저원소 $\eta$를 다른 임의의 기저의 기저원소 $\eta^*$로 바꾸면, 위에서 보았듯이 $\mathfrak G$와 동형인 영역이 생긴다. 역으로 $\mathfrak G^*$를 임의의 영역, 따라서 $\mathfrak L(\mathfrak G)$의 어떤 영역 $\mathfrak G$와 동형인 영역이라 하자. 이때 기저원소 $\eta$에 $\mathfrak G^*$의 원소 $\eta^*$가 대응한다면, 이 원소들은 다시 초월기저를 이룬다. 또한 $\mathfrak G^*$는 $\eta^*$의 대수함수 가운데 $\mathfrak G$가 $\eta$에 관하여 포함하는 것과 같은 함수들, 또는 그 켤레들을 포함한다. \emph{따라서 $\mathfrak G$와 동형인 모든 영역은 $\mathfrak G$와 $(H)$에 관한 그 켤레 영역들에서}\srcfn{*)}{$\mathfrak G$의 켤레 영역에 속하는 원소들은 II에 따라 분명 $\mathfrak G$의 원소들을 써서 유리적으로 나타낼 수 있다. 그러나 반드시 다항식적으로 나타낼 수 있는 것은 아니므로, 켤레 영역은 $\mathfrak G$와 다를 수 있다.} \emph{--- 켤레 영역이 $\mathfrak G$와 다를 경우에는 이들도 포함하여 --- 초월기저 $H$가 가능한 모든 초월기저를 두루 취하게 함으로써 얻어진다.} $\mathfrak L(\mathfrak G)$의 고정된 각 영역 $\mathfrak G_i$에서 이 방법으로 유별 $\mathfrak R_i$가 생긴다. 이 유별은 $\mathfrak G_i$와 동형인 모든 영역 $\mathfrak G$만을 빠짐없이 포함하지만, 같은 영역을 여러 번, 심지어 무한히 여러 번 포함할 수도 있다. 위의 방법으로 $\mathfrak R_i$에서 $\mathfrak G_i$와 동형인 각 영역을 정확히 한 번씩만 포함하는 부분집합 $\mathfrak R_i^*$를 골라낼 수 있다. $\mathfrak R_i^*$가 모든 유별을 두루 취하게 하면, 모든 영역 $\mathfrak G$를 각각 정확히 한 번씩 얻는다. 물론 유별 $\mathfrak R_i$를 특징짓는 데에는 $\mathfrak G_i$ 대신 그 유별의 다른 어느 영역을 써도 된다. 그 영역에서도 위와 같이 $\mathfrak R_i$의 모든 영역을 이끌어 낼 수 있기 때문이다. 요약하면 다음과 같다.

\emph{$\mathfrak M(\mathfrak G)$에서 부분집합 $\mathfrak L(\mathfrak G)$를 골라내어, 그 영역 $\mathfrak G$들을 모든 유별의 전체와 중복 없이 일대일로 대응시킬 수 있다. 이때 $\mathfrak G_i$에 딸린 유별 $\mathfrak R_i$는 $\mathfrak G_i$와 $(H)$에 관한 그 켤레 영역들에서 $H$가 가능한 모든 초월기저를 두루 취하게 함으로써 생긴다. $\mathfrak R_i^*$가 $\mathfrak R_i$의 서로 다른 모든 영역을 뜻하고, $\mathfrak R_i^*$가 모든 유별을 두루 취하게 하면, 모든 영역 $\mathfrak G$를 각각 정확히 한 번씩 얻는다.}\srcfn{*)}{$\mathfrak M(\mathfrak G)$는 § 6과 § 7에 따라 다음과 같이 구성한다. § 7에 따라 초월기저 $H$를 그것에서 생기는, 부가조건을 만족하는 각 유리기저 $\Theta$로 확장하고, 각 $\Theta$에 대하여 § 6 b)에 따라 $\Theta$를 포함하는 모든 영역 $\mathfrak G$의 전체 $\mathfrak N(\mathfrak G)$를 만든다.}

대수적으로 정수인 초월수로 이루어진 영역 $\mathfrak H$에 대해서도 같은 사실이 성립한다.

\section*{§ 9. 영역 $\mathfrak G$의 가산무한히 많은 유별의 한 예.}

§ 4 b)와 § 6 b)에 제시한 구성들의 기수는 그곳의 주석에서 알 수 있듯이 모든 정렬순서의 기수와 같고, 따라서 연속체의 기수를 $c$라 할 때 $2^c$와 같다. 그러나 여기서 각 영역을 정확히 한 번씩 센 모든 영역 $\mathfrak G$의 전체가 갖는 기수를 역으로 알아낼 수는 없으며, 모든 유별의 기수는 더더욱 알아낼 수 없다. 유별의 수가 유한하지 않다는 사실을 § 5의 영역들과 비슷한 가산무한히 많은 영역 $\mathfrak G$의 예로 보이겠다. 이 영역들은 모두 본질적으로 서로 다름이 드러날 것이므로 가산무한히 많은 유별을 준다.\srcfn{**)}{§ 6의 영역 $\mathfrak G$들도 마찬가지로 가산무한히 많은 유별을 주지만, 그 증명은 조금 더 복잡하다.}

§ 5의 (1)과 마찬가지로 영역 $\mathfrak G_\sigma$가 다음 모든 양으로 이루어진다고 하자.
\[
\begin{array}{@{}l@{\quad}l@{}}
(1)& \displaystyle z\equiv a_0+a_1(\eta)+\cdots+a_{\sigma-1}(\eta)
 \pmod{\mathfrak M_\sigma(\eta)}.
\end{array}
\]
\emph{여기서 $a_i(\eta)$는 $\eta$의 $i$차 동차형식이고, 모듈 $\mathfrak M_\sigma$는 $\eta$의 $\sigma$차 멱단항식 전부를 기저로, $\eta$의 모든 정수적 대수함수를 승수로 가진다.}\srcfn{***)}{물론 개별 $z$마다 모듈에 나타나는 $\eta$의 멱단항식은 유한개뿐이다.} 두 영역 $\mathfrak G_\sigma$와 $\mathfrak G_\tau$ $(\sigma\ne\tau)$ 사이에 동형 대응이 있으면 모듈 $\mathfrak M_\sigma$와 $\mathfrak M_\tau$도 반드시 서로 동형으로 대응함을 보이겠다. 그러면 그러한 동형 대응이 불가능하다는 결론이 쉽게 따른다.

$\mathfrak G_\tau$의 부정원들을 $\xi$라 하고, 동형 대응에 따라 이 $\xi$들에 $\mathfrak G_\sigma$의 양 $f(\eta)$들이 대응한다고 하자. 동형 대응은 몫을 취해도 보존되므로, $\xi$의 모든 정수적 대수함수로 이루어진 영역 $\mathfrak G_\xi$에는 대수적 정수성에 관하여 닫힌 영역 $\mathfrak T$가 대응한다. $\mathfrak T$는 $f(\eta)$에 대하여 대수적으로 정수인 모든 함수로 이루어지며, 따라서 $\eta$에 관해서도 대수적으로 정수이다.

(1)의 한 양에 $\mathfrak M_\tau$의 한 양이 대응한다고 하자. $\mathfrak M_\tau$의 모듈 성질에 따라 (1)에 $\mathfrak T$의 임의의 양 $\psi(\eta)$를 곱한 것도 $\mathfrak G_\sigma$에 속해야 한다. 식으로 쓰면
\[
\begin{array}{@{}l@{\quad}l@{}}
(2)& \displaystyle (a_0+a_1(\eta)+\cdots+a_{\sigma-1}(\eta))\psi(\eta)
\equiv b_0+b_1(\eta)+\cdots+b_{\sigma-1}(\eta)
\pmod{\mathfrak M_\sigma}.
\end{array}
\]
모든 $\eta$를 영으로 놓으면 $a_0\psi(0)=b_0$를 얻으므로, (2)는 다음과 같이 바뀐다.
\[
\begin{array}{@{}l@{\quad}l@{}}
(3)& \displaystyle (a_0+a_1(\eta)+\cdots+a_{\sigma-1}(\eta))(\psi(\eta)-\psi(0))
\equiv c_1(\eta)+\cdots+c_{\sigma-1}(\eta)
\pmod{\mathfrak M_\sigma}.
\end{array}
\]

$\psi(\eta)$와 더불어 각 $\nu$에 대하여 함수
\[
 \chi_\nu(\eta)=\sqrt[\nu]{\psi(\eta)-\psi(0)}
\]
도 $\mathfrak T$에 속하고 $\chi_\nu(0)=0$이므로, (3)과 마찬가지로 다음 식들이 성립한다.
\[
\begin{array}{@{}l@{\quad}l@{}}
(4)& \displaystyle (a_0+a_1(\eta)+\cdots+a_{\sigma-1}(\eta))\chi_\nu(\eta)
\equiv c_1^{(\nu)}(\eta)+\cdots+c_{\sigma-1}^{(\nu)}(\eta)
\;\mathrm{mod.}\;\mathfrak M_\sigma
\quad(\nu=1,2,\ldots).
\end{array}
\]
$\chi_1=\psi(\eta)-\psi(0)$와 그 $\kappa-1$개 켤레의 곱을 차수 $\lambda$인 $A(\eta)$라 하자. $\chi_1(0)=0$이므로 $A(\eta)$는 차수가 적어도 1인 항들로 시작해야 한다. (4)에서
\[
\begin{array}{@{}l@{\quad}l@{}}
(5)& \displaystyle a_0\cdot\chi_\nu(\eta)\equiv0\;\mathrm{mod.}\;\mathfrak M_1;\qquad
a_0^\nu\cdot\chi_1(\eta)\equiv0\;\mathrm{mod.}\;\mathfrak M_\nu;\qquad
a_0^{\nu\kappa}\cdot A(\eta)\equiv0\;\mathrm{mod.}\;\mathfrak M_{\nu\kappa}
\end{array}
\]
를 얻는다. 따라서 $a_0\ne0$이면 § 3에서와 마찬가지로 $\lambda>\nu\kappa$이다. 그러나 $\nu<\frac{\lambda}{\kappa}$라는 제한은 $\mathfrak T$가 대수적 정수성에 관하여 닫혀 있다는 사실과 모순되므로, 반드시 $a_0=0$이다. 이제 (4)에서
\[
\begin{array}{@{}l@{\quad}l@{}}
(6)& \displaystyle a_1(\eta)\cdot\chi_\nu(\eta)\equiv c_1^{(\nu)}(\eta)\;\mathrm{mod.}\;\mathfrak M_2;\qquad
[a_1(\eta)]^{\nu\kappa}\cdot A(\eta)
\equiv[c_1^{(\nu)}(\eta)]^{\nu\kappa}\;\mathrm{mod.}\;\mathfrak M_{\nu\kappa+1}
\end{array}
\]
을 얻는다. $A(\eta)$가 차수 적어도 1인 항들로 시작하므로 여기서 $c_1^{(\nu)}(\eta)=0$이 따른다. 그러므로 (5)와 완전히 마찬가지로 각 $\nu$에 대하여
\[
 a_1(\eta)\cdot\chi_\nu(\eta)\equiv0\;\mathrm{mod.}\;\mathfrak M_2;\qquad
 [a_1(\eta)]^{\nu\kappa}\cdot A(\eta)\equiv0\;\mathrm{mod.}\;\mathfrak M_{2\nu\kappa}
\]
이고, 여기서 $a_1(\eta)=0$을 얻는다. 이 과정을 계속하여 (5)와 (6)을 번갈아 적용하면
\[
 a_0=0,\qquad a_1(\eta)=0,\quad \ldots,\quad a_{\sigma-1}(\eta)=0
\]
을 얻는다. $\mathfrak G_\tau$에 관해서도 같은 논법이 성립하므로, \emph{$\mathfrak G_\sigma$와 $\mathfrak G_\tau$ $(\sigma\ne\tau)$ 사이의 모든 동형 대응에서 모듈 $\mathfrak M_\sigma$와 $\mathfrak M_\tau$는 서로 동형으로 대응한다.}

이를테면 $\sigma>\tau$라고 하자. 부정원 $\xi$에는 $\mathfrak G_\sigma$의 양 $f(\eta)$가 대응했다. $\mathfrak G_\tau$의 모든 양은 $\mathfrak M_\tau$를 법으로 하여 $\xi$의 다항식으로 나타낼 수 있고, $\mathfrak M_\tau$와 $\mathfrak M_\sigma$가 서로 대응하므로 $\mathfrak G_\sigma$의 모든 양은 $\mathfrak M_\sigma$를 법으로 하여 $f(\eta)$의 다항식으로 나타난다. $\sigma>\tau\ge1$이므로 부정원 $\eta$는 분명 $\mathfrak M_\sigma$에 속하지 않는다. 따라서 $\eta$를 $\mathfrak M_\sigma$를 법으로 하여 $f(\eta)$의 다항식으로 나타낼 수 있으려면, 선형항이 사라지지 않는 $f(\eta)$가 있어야 한다. 그러한 대응을
\[
 \xi:f(\eta)\equiv a_0+a_1(\eta)\pmod{\mathfrak M_2}\quad[a_1(\eta)\ne0],
\]
\[
 \xi^\tau:[f(\eta)]^\tau\equiv a_0^\tau\pmod{\mathfrak M_1}\quad\text{$a_0\ne0$일 때},
\]
\[
 \hphantom{\xi^\tau:[f(\eta)]^\tau}\equiv [a_1(\eta)]^\tau
 \pmod{\mathfrak M_{\tau+1}}\quad\text{$a_0=0$일 때}
\]
라고 쓰자. $\xi^\tau$는 $\mathfrak M_\tau$에 속한다. 그러나 $[f(\eta)]^\tau$는 사라지지 않는 0차항 또는 $\tau$차항으로 시작하고 $\tau<\sigma$이므로 $\mathfrak M_\sigma$에 속할 수 없다. 이는 앞서 얻은 $\mathfrak M_\sigma$와 $\mathfrak M_\tau$의 대응과 모순이다. $\tau>\sigma$인 경우는 $\eta$와 $\xi$를 맞바꾸면 함께 처리되므로, 영역 $\mathfrak G_\sigma$와 $\mathfrak G_\tau$가 \emph{서로 동형이 아니며, 곧 본질적으로 서로 다름}이 증명되었다. \emph{따라서 영역 $\mathfrak G_i$ $(i=1,2,\ldots)$는 영역 $\mathfrak G$의 가산무한히 많은 유별의 한 예를 준다.}

\emph{주.} 이처럼 서로 다른 두 영역 $\mathfrak G_i$는 동형이 아니지만, 두 영역 각각이 다른 영역의 한 부분정역과 동형일 수는 있다.\srcfn{*)}{이러한 현상은 체에서도 나타날 수 있다. Steinitz, 앞의 논문 끝부분 참조.} 이를테면 $\tau=1$이고 $\sigma$는 임의라고 하자.
\[
\mathfrak G_1:a_0+\mathfrak M_1(\xi);\qquad
\mathfrak G_\sigma:a_0+a_1(\eta)+\cdots+a_{\sigma-1}(\eta)+\mathfrak M_\sigma(\eta).
\]
대응 $\xi=\eta$에 따라 $\mathfrak G_\sigma$는 $\mathfrak G_1$의 진부분정역이 된다. 역으로 $\mathfrak G_1$과 $\mathfrak G_\sigma$에서 일대일인 대응 $\xi=\eta^\sigma$, $\eta=\sqrt[\sigma]{\xi}$에 따라 $\mathfrak G_1$은 $\mathfrak G_\sigma$의 진부분정역이 된다. 따라서 각 $\mathfrak G_i$는 자기 자신의 진부분정역과도 동형이다. 그러므로 교집합이라는 개념과 달리 교집합 유별, 곧 그 유별의 각 영역이 다른 모든 유별의 각 영역의 진부분정역과 동형인 그러한 유별은 존재하지 않는다. 적어도 유일하게 정해지는 것으로는 존재하지 않는다.

\section*{§ 10. 임의의 체의 정수적 양.}

앞 절들에서 전개한 내용은, 영역 $\mathfrak G$의 정의 성질 III과 IV에 대수적 수를 포함시키지 않는 한, $\mathfrak G$에 관해서는 모든 복소수의 전체가 체를 이룬다는 사실에만 의존했다.\srcfn{**)}{§ 7에서만 유리수의 ``소체''가 무한히 많은 원소를 갖는다는 추가 가정을 사용했다. 원소가 유한개뿐인 소체의 경우에는, 뒤에서 보듯이 이 절이 단순히 빠진다.} 영역 $\mathfrak H$에 관해서는 이 체가 대수적으로 닫혀 있다는 사실을 추가로 사용했다. 따라서 얻은 결과들은 Weber와 E. Steinitz가\srcfn{*)}{앞의 논문 서론.} ``결합법칙과 교환법칙을 따르고 분배법칙으로 연결되며, 제한 없이 유일하게 역연산을 할 수 있는 두 연산, 곧 덧셈과 곱셈을 갖춘 원소들의 계''라고 추상적으로 정의한 \emph{가장 일반적인 체}로 곧바로 일반화된다.\srcfn{**)}{영으로 나누는 것만은 제외해야 한다.}

그러한 모든 체는 단위원을 포함하고, 따라서 단위원에서 덧셈과 곱셈 및 그 역연산으로 이끌어 낸 ``소체''를 포함한다. 이 소체는 유리수형(표수 0)이거나 소수 $p$에 대한 잉여류계형(표수 $p$)이다.\srcfn{***)}{Steinitz, § 4.} 이때 \emph{소체의 정수적 양}은 단위원에서 덧셈과 뺄셈으로 생기는 양으로 잘 정의된다. 표수 $p$의 소체에서는 정수적 양이 이미 소체 전체를 소진하므로 소체의 분수형 양은 존재하지 않는다. 모든 대수적으로 닫힌 체는 소체에 관하여 대수적인 모든 원소를 소체에 첨가하여 얻는 ``절대대수체''를 포함한다. 따라서 표수 0인 체에서는 이 체가 모든 대수적 수의 체와 같은 유형이다. \emph{절대대수체의 대수적 정수인 양}은 소체의 정수적 양에 대하여, 또는 같은 말로 단위원에 대하여 대수적으로 정수인 모든 양으로 잘 정의된다. 그러므로 표수 $p$의 체에서는 절대대수체의 분수형 대수적 양도 존재하지 않는다. 이러한 예비 고찰에 따라 ``정수적'' 양과 ``대수적으로 정수인'' 양의 정의를 각각 임의의 체와 임의의 대수적으로 닫힌 체로 옮길 수 있다. \emph{체 $\Omega$의 정수적 양은 $\Omega$의 양들로 이루어진 정역 $\mathfrak G$로 정의한다. 이 정역의 분수체}\srcfn{$\dagger$)}{즉 $\mathfrak G$의 두 양의 모든 몫으로 이루어진 체.} \emph{는 $\Omega$ 전체와 같고, 이 정역은 단위원을 포함하지만 소체의 분수형 양은 포함하지 않는다.}

\emph{대수적으로 닫힌 체 $\Omega$의 대수적 정수인 양은 $\Omega$의 양들로 이루어지고 대수적 정수성에 관하여 닫힌 정역 $\mathfrak H$로 정의한다. 이 정역의 분수체는 $\Omega$ 전체와 같고, 이 정역은 단위원을 포함하지만 절대대수체의 분수형 대수적 양은 포함하지 않는다.}

표수 0인 체에 대해서는 앞 절들에서 얻은 결과들이 글자 그대로 성립한다. 다만 ``대수적 수''를 각각 소체 또는 절대대수체의 양으로 바꾸면 된다. 표수 $p$인 체에서는 소체에 분수형 양이 없기 때문에 결과들이 훨씬 더 단순해진다. 유리기저와 첨가할 계에 관한 부가조건이 완전히 사라지는 것이다.\srcfn{*)}{이 절의 첫 주석 참조.} 따라서 결과는 다음과 같다. \emph{추상적 정의에 따라 표수 $p$인 체의 정수적 양으로 허용되는 것은 임의의 유리기저에서 생기는 각 정역, 체 자체, 그리고 이 정역들과 체 사이에 놓인 모든 정역이다. 이에 대응하여 소수 표수의 대수적으로 닫힌 체에서 대수적 정수인 양으로 간주되는 것은, 초월기저에서 생긴 정역과 체 자체 사이에 놓인 대수적 정수성에 관하여 닫힌 각 정역이며, 두 경계 영역도 포함된다.} 따라서 소수 표수의 가장 일반적인 체에서는 그 소체에서와 마찬가지로 정수적인 것과 분수형인 것의 차이가 희미해진다.

\medskip
\noindent 에를랑겐, 1915년 3월 30일.
